\section{From layout cost to service cost}
\label{sec:mappingcosts}
\subsection{Mapping execution and representation costs}
\subsubsection{Execution mode changes the comparison}
We compare PQ, QB and query-to-head mapping (QH), which places independent queries in a head-like dimension and then restores native output order. Each operator study freezes its inputs, reference and cost window~\cite{mappingrecords,b570compactrecords}.

Graph allocations initially reused K/V storage and changed later operands. Independent persistent storage restored repeatable input and output bytes. Balanced experiments then used identical QB arms to measure A/A variation. On the RTX core panel, QB\_A/QH is 0.900826. A frozen selector predicts the ordering on 12 held-out synthetic geometries, but charging its lookup raises the offline estimated selector/fixed-QB ratio to 1.010298. Selection overhead outweighs the predicted benefit.

CUDA Graph capture, instantiation and replay materially change the relative costs. QB\_A's completed-graph ON/OFF ratios range from 0.615 to 0.849. For four permuted-$q=2$ coordinates, the QH-minus-QB gap falls from 12.056--12.975 to 0.735--1.627\,$\mu$s. Instrumented mode checks and clean timing runs share frozen graph/runtime identities. This intervention identifies an execution-mode effect on QB versus QH; it does not measure the PQ/QB mechanism in the resident-request study.

\subsubsection{Path availability and reference preservation}
RTX key extents 1023/1024/1025 select 8/3/8-kernel paths, with completed graph-replay costs of 113.085/50.439/114.618\,$\mu$s. To separate physical extent from logical visibility, a padding intervention fixes each input and original mask columns, then appends finite zero K/V and fully masked columns~\cite{mappingrecords}.

Padding logical length 1023 to physical extent 1024 reaches the three-kernel path on the first call but fails the unpadded reference on all four qualification hosts. Padding to 1025 passes all four, retains eight kernels and costs about 1.5\% more than either native QB control. The faster path is therefore available but inadmissible under the frozen reference (Figure~\ref{fig:mappingboundary}); it has no qualified steady-state cost.

\begin{figure*}[t]
\centering\begin{tikzpicture}[font=\footnotesize]
\node[anchor=north west,inner sep=0] at (0,0) {\begin{minipage}[t]{.48\textwidth}
\textbf{A\quad Fixed logical length 1023: RTX}\par\vspace{5pt}
\begin{tabular}{@{}rrrr@{}}\toprule
Physical extent & Qualification & Kernels & Cost / native A\\\midrule
1023 (native) & 4/4 pass & 8 & 1.000000\\
1024 (padded) & \textcolor{draft}{0/4 fail} & 3$^\ast$ & ---\\
1025 (padded) & \textcolor{qualified}{4/4 pass} & 8 & 1.014783\\\bottomrule
\end{tabular}\par\vspace{5pt}
$^\ast$First-call trace only; the failed path has no qualified clean cost. Logical inputs and original mask columns are unchanged.
\end{minipage}};
\node[anchor=north west,inner sep=0] at (.515\textwidth,0) {\begin{minipage}[t]{.48\textwidth}
\textbf{B\quad Complete Q-representation transactions}\par\vspace{5pt}
\begin{tabular}{@{}llr@{}}\toprule
Panel & Comparator & Compact-QH ratio\\\midrule
RTX $q=2$ & QB\_A / QB\_B & 1.002216 / 1.008295\\
RTX $q=4$ & QB\_A / QB\_B & 1.001598 / 1.004167\\
B570 $q=2,4$ & dual-QB & 0.992244\\\bottomrule
\end{tabular}\par\vspace{5pt}
B570 adjusted interval: [0.979908, 1.004691]. All tested practical-effect screens remain unresolved; no QH advantage is established.
\end{minipage}};
\end{tikzpicture}

\caption{Reference and transaction constraints on mapping optimization. A: the length-1023 padding candidate reaches a three-kernel path on its first call but fails the reference. B: complete Q-representation transactions with explicit comparators. RTX and B570 use different geometries and execution modes. The B570 interval is descriptive and family-adjusted.}
\label{fig:mappingboundary}
\end{figure*}

\subsubsection{Charging the representation transaction}
Compact Q removes a QH graph copy but must first be produced. The RTX timer covers the continuous transaction from common permuted host Q through copy/transpose, upload, graph replay/synchronization and complete ordinary-output readback. K/V and mask are resident; targets are preallocated. The window excludes model Q production and ongoing KV updates.

For $q=2$ and $q=4$, compact-QH/QB\_A transaction ratios are 1.002216 and 1.001598; relative to QB\_B they are 1.008295 and 1.004167. Neither passes the practical screen against both controls. The B570 equal-head study likewise establishes no practical QH benefit: compact-QH/dual-QB is 0.992244, with descriptive adjusted interval [0.979908,1.004691] and paired difference $-1.518\,\mu$s [$-4.168,+0.916$]. \sref{supp:followupmapping} retains the separate resident-graph results.

Fixed QB remains the baseline in these domains: a representation must repay its production and consumption costs, not just reduce internal graph work.
